\section{Introduction}
\label{sec:intro}

Connectomes are now available at a scale that invites a structural explanation of
neural computation: the complete somatic wiring of \textit{C. elegans}
\cite{cook2019whole}, the adult \textit{Drosophila} brain \cite{dorkenwald2024neuronal,
lin2024network}, proofread mammalian cortical volumes \cite{microns2025functional},
and whole-brain human networks \cite{vanessen2012human}. Alongside them has grown a
family of claims of the following shape: biological wiring contains statistically
overrepresented subgraphs -- feedforward loops, reciprocal inhibitory pairs -- these
motifs implement identifiable computations, and networks built to contain them
inherit an advantage in accuracy, robustness or energy over conventional artificial
networks.

Each step of that chain is plausible, and the first step is well established in the
sense that motif overrepresentation against degree-preserving nulls is a reproducible
empirical finding \cite{milo2002network,milo2004superfamilies,sporns2004motifs} --- our
own recomputation finds it too, though less uniformly across species than the phrase
``conserved motifs'' suggests (\S\ref{sec:res-motifs}). The later steps are where the
evidence is thinner than the claims, for a reason that is methodological rather than
empirical.

\paragraph{A contradiction, and what it turned out to be about.}
The sharpest quantitative support for the later steps comes from reservoir computing, and
it is the part of this literature that reaches furthest: it is routinely offered as a
reason to build neuromorphic systems on connectome-derived wiring.
Su\'arez et al.~\cite{suarez2024connectome} report that empirical connectomes have
significantly greater linear memory capacity than degree-preserving rewirings. Our own
ladder, run in a spiking network at a higher gain, found no such advantage.

We set out to adjudicate that disagreement and did not manage to. What we found instead
is that the comparison ``does a connectome beat its rewiring?'' is not one experiment. It
is a composition of four choices, each of which is normally reported in a subordinate
clause, and each of which can change the sign of the answer on identical data.

\paragraph{Computational coupling.}
We call the quantity this literature is trying to measure the \emph{computational
coupling} of a network: the extent to which what a network computes is attributable to
how it is wired. Coupling is never observed directly. What is observed is a score on a
task, and a score belongs to the network together with everything the measurement
supplies around it --- the units the network is built from, the regime it is driven in,
the channels it is read through, and the alternative networks it is scored against. The
coupling is the part of that score which moves when the wiring moves and nothing else
does. It is therefore defined only relative to a stated procedure for moving the wiring,
and it is the sign of that difference, not the score itself, that the structural claims
in this literature rest on.

\begin{definition}[Computational coupling]
\label{def:coupling}
Let $C_{U,\rho,R}(W)$ denote the task capacity of a network with weight matrix $W$,
evaluated with unit model $U$, at operating point $\rho$, under readout protocol $R$;
and let $\mathcal{N}$ be an ensemble of surrogate networks that alters the wiring while
holding stated properties of $W$ fixed. The \emph{computational coupling} of $W$ under
the measurement protocol $P = (\mathcal{N}, U, \rho, R)$ is the standardised contrast
\begin{equation}
\label{eq:coupling}
\kappa(W; P) \;=\;
\frac{C_{U,\rho,R}(W) \;-\; \mathbb{E}_{W' \sim \mathcal{N}}\, C_{U,\rho,R}(W')}
     {\operatorname{sd}_{W' \sim \mathcal{N}}\, C_{U,\rho,R}(W')}\,.
\end{equation}
\end{definition}

\noindent
Equation~\ref{eq:coupling} is not a new statistic: it is the $Z$-score the field already
reports, under connectome-versus-null comparisons, motif enrichment, small-worldness and
rich-club analysis alike. Naming it is worth the trouble for one reason --- the name
carries its arguments with it. ``The connectome beats its rewiring'' is the claim
$\kappa > 0$, and $\kappa$ is a function of four things besides $W$. Each of
$\mathcal{N}$, $U$, $\rho$ and $R$ is a choice; each is normally reported in a
subordinate clause, if at all; and each, we find, can change the sign of $\kappa$ on
identical data. A coupling reported without its protocol is not a measurement of the
connectome, because the same connectome supports both signs. We measured all four
choices (Table~\ref{tab:four-choices}).

\begin{enumerate}[leftmargin=1.4em,itemsep=3pt,topsep=4pt]
\item\textbf{Which rewiring counts as the null} ($\mathcal{N}$)\textbf{.} ``Degree-preserving randomisation'' names
at least two different ensembles, depending on whether a weight travels with its edge or
is re-assigned afterwards. Both match the phrase; they are not the same distribution over
graphs; and on the same connectomes, at the same operating point, in the same unit model,
they do not agree about the sign of the result (\S\ref{sec:res-nullchoice}).

\item\textbf{What a unit is} ($U$)\textbf{.} A $\tanh$ unit is never silent; a sub-threshold
integrate-and-fire unit contributes nothing at all. Running identical graphs, identical
rewirings, identical input and an identical readout through both --- on the five
connectomes distributed with the original toolbox --- changes the outcome, and it matters
because the hardware the result is invoked for is spiking (\S\ref{sec:res-unit}).

\item\textbf{Where on the gain axis it is measured} ($\rho$)\textbf{.} The published comparison is located
\emph{at criticality}, at a spectral radius of $1$, the regime a reservoir is
conventionally evaluated in \cite{bertschinger2004real,legenstein2007edge,%
boedecker2012information}. The connectome-versus-null contrast changes sign along that
axis --- in $\NGainSignFlip$ of $\NGainNetworks$ networks, significant on both sides --- so
a comparison reported at one gain is a statement about that gain
(\S\ref{sec:res-operating}).

\item\textbf{How the reservoir is read out} ($R$)\textbf{.} Under a protocol in which every neuron is
both driven and decoded, a reservoir with its recurrent matrix \emph{set to zero} scores
$\ZeroRecFull$ against $\ZeroRecFullConn$ for the connectome. In that regime the
comparison is not measuring the wiring at all (\S\ref{sec:res-protocol}).
\end{enumerate}

\begin{table*}[t]
\centering\small
\caption{\textbf{Four choices that a connectome-versus-null comparison has to make, and
what each one is worth.} Each row is a decision normally reported in a subordinate clause
of a methods section. For each we ran both settings on otherwise identical data and
report the largest effect the choice alone produces. None of the four is a robustness
check in our design: each can change the sign of the reported result, and the first and
fourth can change it between two significant verdicts in opposite directions.}
\label{tab:four-choices}
\begin{tabular}{p{0.15\textwidth}p{0.25\textwidth}p{0.25\textwidth}p{0.25\textwidth}}
\toprule
\textbf{Choice} & \textbf{Setting A} & \textbf{Setting B} & \textbf{What the choice alone
does} \\
\midrule
Null ensemble \newline (\S\ref{sec:res-nullchoice})
& BCT \texttt{randmio}: weights travel with their edges
& Rank-reassigned: same weight multiset, re-placed by degree rank
& Opposite signs in $\NNullSignFlip$ of $\NNullCells$ cells; opposite significant
verdicts in $\NNullVerdictFlip$; median disagreement $\NullDisagreeMedian$ capacity
units \\
\addlinespace
Unit model \newline (\S\ref{sec:res-unit})
& Graded $\tanh$, never silent
& Leaky integrate-and-fire, silent below threshold
& Capacity compressed $\UnitCapRatio\times$ for every topology alike; the contrast
survives in $\NUnitSpkLeads$ of $\NUnitSpkNetworks$ networks \\
\addlinespace
Operating point \newline (\S\ref{sec:res-operating})
& At criticality, $\rho = 1$
& Above it, $\rho > 1$
& Contrast changes sign along the axis in $\NGainSignFlip$ of $\NGainNetworks$ networks,
significant on both sides \\
\addlinespace
Readout protocol \newline (\S\ref{sec:res-protocol})
& All-to-all: every neuron driven and decoded
& Sensory/motor-disjoint
& Under A, zeroing the recurrent matrix entirely costs nothing ($\ZeroRecFull$ against
$\ZeroRecFullConn$); under B it falls to chance ($\ZeroRecSplit$) \\
\bottomrule
\end{tabular}
\end{table*}

This is the paper's central claim, and it is a claim about how the field measures rather
than about what brains are. Computational coupling is a property of a connectome
\emph{and} the apparatus that measures it, and the four arguments of $\kappa$ are that
apparatus. Each of the four is a fact about the measurement that is indistinguishable,
in a single reported number, from a fact about the topology. The fourth is the limiting
case: under an all-to-all readout the coupling is driven to zero by construction, since a
reservoir whose recurrent matrix has been deleted outright scores as well as the
connectome, and a comparison run there reports a number that no change to $W$ could have
moved. Our own
results carry exactly the same dependence, which is why we report the axis and not the
point for them too --- and why, where our design could not resolve an effect, we say so
rather than reporting a non-detection as a refutation.

\paragraph{The attribution problem.}
Establishing that a structural property \emph{causes} a computational advantage
requires a control that differs in that property and nothing else. The standard
comparison -- a connectome against a degree-preserving rewiring, or against an
Erd\H{o}s--R\'enyi graph -- does not have this form. Rewiring a connectome
simultaneously changes its motif census, its clustering, its spectral structure and
every higher-order statistic that co-varies with them, so a performance difference
cannot be attributed to motifs in particular. The ensemble that would settle the question --- one that randomises a connectome while
\emph{holding its motif census fixed} --- is standard in statistical network science and
absent from this debate. Motif-constrained ensembles have been sampled by annealing
\cite{milo2002network}, by exponential random graph models
\cite{frank1986markov,robins2007introduction}, by the $dk$-series
\cite{mahadevan2006systematic,orsini2015quantifying}, and, with a sampling guarantee, by
Wang--Landau multicanonical methods \cite{fischer2015sampling}. None has been used to ask
what a connectome's motifs \emph{do}. We build a sign-attributed version at connectome
scale and place it in a nested ladder of controls in which each rung adds exactly one
structural property, so that the contribution of each can be read off as a difference
between adjacent rungs.

Three further confounds recur in this literature and are worth stating plainly, since
all three appear in our own earlier work on this project. First, a null model described as
preserving degrees \emph{and} Dale's principle preserves nothing beyond degrees: when
synapses inherit their presynaptic neuron's valence, Dale's principle is automatically
invariant under double-edge swaps, so such a null cannot test whether the placement of
inhibition carries information. Second, and more consequential, a reservoir in which
every neuron is both driven and decoded does not require its recurrent connectivity at
all -- we show that zeroing the recurrent matrix entirely leaves SHD accuracy unchanged
under that protocol. A structural comparison run in that regime measures nothing,
whatever it reports. Third, matching spectral radius across networks equalises recurrent
gain but not the activity level each network settles into, and since firing rate both
predicts accuracy and determines the synaptic-operation count outright, a comparison at
matched gain confounds topology with operating point -- in the direction that flatters
the unstructured control. We therefore report the ladder both at matched gain and with
every network recalibrated to a common firing rate.

A common thread runs through all three: each is a failure that produces no error and no
warning. The naive protocol returns plausible accuracies, the Dale-compliant null returns
plausible $Z$-scores, and the matched-gain comparison returns a plausible ordering. In the
first two the object of study turns out to have made no difference to the number reported;
in the third it made one, but not for the reason the comparison was set up to attribute it
to. Nothing in any of their outputs says so. That is why we regard the controls that
expose them as part of the method rather than as robustness checks.

\paragraph{The degeneracy problem.}
A second strand asks the converse question: how much can a circuit be rewired before
function degrades? Biological circuits are famously degenerate
\cite{prinz2004similar,marder2011variability,edelman2001degeneracy}, and the natural
formal tool is the Davis--Kahan theorem \cite{daviskahan1970rotation}, bounding
invariant-subspace rotation by the perturbation size over the spectral gap. Applied to
measured connectomes the bound is vacuous, because their spectral gaps are tiny
($\MinSpectralGap$ at the smallest): it permits a perturbation of norm at most
$\MaxDKAdmissible$ while the networks tolerate a great deal more.
Reporting the resulting discrepancy as a finding about biology inverts the situation --
it is a statement about the instability of eigenvectors under near-degenerate spectra.
We replace it with a bound on the object that task performance actually depends on, the
finite-horizon input--output map, which requires no gap assumption; and we extract from
it a quantity, the operational alignment $\gamma$ of a perturbation, that turns
``biological rewiring is non-adversarial'' into a measurable prediction. Crucially, the
prediction is falsifiable, and it is partly falsified: we construct the maximally
damaging perturbation of matched cardinality and find that $\gamma$ ranks the
perturbation directions by tolerance as predicted, but does not reproduce the $1/\gamma$
scaling the bound implies (\S\ref{sec:res-degeneracy}). We report the failure alongside
the part that holds. What survives is enough for the point at issue: a degeneracy result
demonstrated only against random pruning shows that random directions are benign, which
is a much weaker statement than it appears.

\paragraph{Contributions.}
\begin{enumerate}[leftmargin=1.2em,itemsep=2pt]
\item An exact sign-attributed motif census with $O(d)$ incremental updates, resting on
the observation that under Dale's principle a feedforward loop's coherence is decided
solely by the valence of its intermediate neuron (Eq.~\ref{eq:coherence}).
\item A \textbf{motif-preserving null ensemble}: degree-, valence- and census-matched
surrogates sampled by an annealed swap chain. On \textit{C. elegans} it matches all six
targeted motif counts to within $\CeleAnnealErr$ while retaining only
$\CeleOverlap$ edge overlap with the empirical graph, against $\CeleOverlapLtwo$
for unconstrained rewiring, and it converges on every network we test --- including
diffusion tractography, where almost every edge is reciprocal and an ordinary swap
chain cannot hold the census at all.
\item A \textbf{sampler-robustness analysis} of that ensemble: the same surrogates drawn
over a factor-of-eight range of chain lengths and from both possible starting points
(the empirical graph under a census constraint, and a fully rewired graph annealed
towards the census), which rules out the two ways an annealed sampler could manufacture
the null result reported below.
\item A \textbf{structural decomposition} of reservoir performance on a real
spike-based benchmark, isolating the contribution of density, degree sequence, valence
placement, motif census and residual wiring, with matched controls, effect sizes and
multiplicity correction.
\item A \textbf{null-ensemble sensitivity analysis}: two rewiring constructions that
both satisfy the phrase ``degree-preserving randomisation'' --- one carrying weights with
their edges, one re-assigning them afterwards --- run over the full design, giving the
connectome's advantage opposite significant verdicts in $\NNullVerdictFlip$ of
$\NNullCells$ cells. We built the second believing it was the first, and report the
episode as the worked example it is rather than removing it.
\item A \textbf{unit-model comparison} on the connectomes distributed with
\texttt{conn2res}: identical graphs, rewirings, inputs, delays and readout run through
graded and spiking units, with capacity compressed $\UnitCapRatio\times$ between them, and
with the silencing hypothesis we proposed for the difference supported against the
published null after we had reported it as refuted against the other one.
\item A \textbf{sweep of the gain axis} showing that the connectome-versus-null contrast
changes sign along it, and that against an unstructured graph it does not appear at any
gain we tested.
\item A \textbf{gap-free perturbation bound} on the transfer map, the operational
alignment $\gamma$, and an adversarial-pruning test of the resulting prediction, which
it passes ordinally and fails quantitatively.
\item A \textbf{reproducible open-PDK synthesis} of the accompanying RTL, reporting
only quantities the flow can establish, and withdrawing previously reported commercial
process figures that were never measured.
\end{enumerate}

We regard the fifth item, and the corrections it entails, as part of the contribution
rather than an appendix to it. The audit record accompanying this work documents each
withdrawn claim and the evidence against it.
